\chapter{情動インターラプションと身体表象 $\lambda, \Lambda$}
ナラティブの確信 $\Omega$ の確定は、純粋な大脳皮質内の抽象的計算にとどまらない。確定した解の種類に応じて、大脳辺縁系および自律神経系への直接的な介入が発生する。

本稿では、身体運動および微小表象を以下のように定式化する。
\begin{itemize}
  \item $\lambda$（小文字）：大脳辺縁系の強い情動的介入を伴わない、日常的・機能的な姿勢、視線移動、習慣的ジェスチャー（微小身体表象）。
  \item $\Lambda$（大文字）：扁桃体を中心とする大脳辺縁系の急激な励起（インターラプション）によって強制惹起された、表情・筋緊張・自律神経反応（情動的身体表象）。
\end{itemize}

情動の種類は、$\Lambda$ の添字によって $\Lambda_{\mathrm{anger}}$（怒り）、$\Lambda_{\mathrm{fear}}$（恐怖）、$\Lambda_{\mathrm{shame}}$（恥）、$\Lambda_{\mathrm{joy}}$（歓喜）のように明示する。

大半の人間の脳内では、外部タスクに没頭していない平常時、DMNが優位に立っている。
そのため、微小な感覚ノイズをトリガーとして即座にエゴの利害空間 $\hat{\Omega}_{\mathrm{ego}}$ が呼び出され、瞬時に $\Omega_{\mathrm{ego}}$ が確定する。
このとき生じるダイナミクスは、以下の連鎖として記述される。

%\begin{equation}
%    \boldsymbol{e}_{\mathrm{seq}} 
%    :\xrightarrow[\hat{\Omega}_{\mathrm{ego}}]{\hat{\Omega}'_{\mathrm{ego}}} \Omega_{\mathrm{ego}} 
%    \Longrightarrow 
%    \begin{cases}
%        \Lambda_{\mathrm{emotion}} & (\text{扁桃体励起による自律神経・身体ループバック}) \\
%        A_{d,\mathrm{ego}} & (\text{言語野インタープリターによる自己正当化作話})
%    \end{cases}
%\end{equation}

\begin{equation}
    \boldsymbol{e}_{\mathrm{seq}} 
    :\xrightarrow[\hat{\Omega}_{\mathrm{ego}}]{\hat{\Omega}'_{\mathrm{ego}}} \Omega_{\mathrm{ego}} 
\xrightarrow{} \langle \Lambda_{\mathrm{emotion}}, A_{d,\mathrm{ego}} \rangle \end{equation}

\begin{itemize}
    \item $\Lambda_{\mathrm{emotion}}$：\textbf{身体・自律神経状態ベクトル（ソマティック・マーカー）} \\
    確信 $\Omega_{\mathrm{ego}}$ の成立に伴い、扁桃体（Amygdala）や前島皮質が励起され、自律神経系・内分泌系を介して全身の生理状態を不可逆的に相転移させる下行性シグナルである。心拍変動、呼吸の狭窄、内臓筋の緊張、血流の偏向といった身体状態の歪み（$\Lambda$）として具現化し、自己の生存に対する利害（profit / threat）を身体レベルで即座に確定させる。

    \item $A_{d,\mathrm{ego}}$：\textbf{言語野インタープリターによる自己正当化作話（内部対話）} \\
    Broca野およびWernicke野をはじめとする言語ネットワークが、非言語的な確信 $\Omega_{\mathrm{ego}}$ と先行して励起した身体状態 $\Lambda_{\mathrm{emotion}}$ の「つじつまを合わせる」ために合成したデジタルな内的ナラティブである。エゴの存在価値やヒエラルキーを保全すべく、過去の記憶から都合の良い証拠をサンプリングし、自己正当化や他者批判の文脈を論理的に捏造する。
\end{itemize}

重要なのは、この両者が独立して存在するのではなく、まず $\Lambda_{\mathrm{emotion}}$ によって身体が情動的にロックされ、その身体感覚を後追いで説明・補強するために $A_{d,\mathrm{ego}}$ が駆動するという非対称な依存関係を持つ点である。




確定した確信 $\Omega_{\mathrm{ego}}$ は、即座に大脳辺縁系へ割り込みをかけて強烈な情動反応 $\Lambda_{\mathrm{emotion}}$（血管収縮、表情の硬直、呼吸停止など）を引き起こす。
同時に、その身体変化そのものが新たな強力なノイズ入力となって脳へループバックし、インタープリターによる独白 $A_{d,\mathrm{ego}}$ を燃料として供給し続けることで、エゴの閉鎖系回路を物理的に固定化（ロックイン）してしまうのである。

リチャード・バンドラーが卓越した洞察力でクライアントの「呼吸の変化」や「瞳孔の散大」「顔面の微小筋緊張」を読み取っていたのは、相手が発する表層の言語スクリプトの背後で、どのナラティブ確信 $\Omega$ が確定し、どの身体インターラプション $\Lambda$ が発火したかをリアルタイムに監視していたからにほかならない。


